\subsection*{E4.Pinching: an excluded route to nonschematic Picard spaces}

\paragraph{Proposition.}
Suppose $X'\to X$ is a finite schematically dominant surjection
of proper $S$-schemes, with its conductor square, and both conductor schemes are finite over
$S$. If $X'\to S$ is locally projective, so is $X\to S$.
Consequently, when $X/S$ is also flat, finitely presented and has
geometrically integral fibres, its full Picard functor is a scheme.

\paragraph{Proof.}
Work near any point $s\in S$ and choose a relatively ample invertible
sheaf $L'$ on $X'$. Its restriction to the inverse-image conductor is
trivial after shrinking $S$ around $s$: over the local scheme
$\operatorname{Spec}\mathcal O_{S,s}$, that conductor is finite and
hence semilocal, so every invertible module is free. The trivialization
spreads to a neighborhood by finite presentation. The conductor
fibre-product square therefore glues $L'$ to the trivial bundle on
the conductor of $X$, producing an invertible sheaf $L$ on $X$.
Since its pullback under the finite surjection is ample, $L$ is
relatively ample. This is the ampleness argument in
\cite[Proposition~2.4(i), p.~6]{BCEBrion2015}; for this argument only
finiteness of the conductors is needed. The final assertion is the
projective Picard theorem.

Thus finite-conductor pinching of a locally projective simultaneous
normalization, in particular of $\mathbf P^1_S$, cannot realize a
nonschematic Picard space. This does not assert that an arbitrary
pinching is flat or that normalization commutes with base change.
The general integral-curve schematicity question remains when no
such normalization over the original base is available. For normal
families of higher dimension, the arbitrary-base question likewise
remains unresolved here.
